\chapter{Évaluation sur 2024 et portée des résultats}
\label{chap:test-ml}

Le notebook 07 évalue la configuration retenue sur 2024. Il ne cherche pas un nouveau
modèle et ne réoptimise ni la représentation des variables, ni $C$, ni le seuil.
Cette dernière étape vérifie si les résultats de validation se maintiennent sur
une année ultérieure.

\section{Reconstruire la configuration avant de calculer le test}
Le fichier exporté par le notebook 06 conserve les variables, les tranches d'âge,
le traitement des absences, $C$, le seuil, les versions et les empreintes des entrées.
Le notebook 07 vérifie ces empreintes puis reconstruit le modèle sur 2019-2022.
L'année 2023 n'est pas ajoutée à l'apprentissage avant le test.

Les scores reconstruits de 2023 sont comparés aux scores enregistrés, avec une
tolérance absolue de $10^{-4}$. Dans les exports fournis, les écarts sont nuls.
Ce contrôle est effectué avant le calcul des performances de 2024. Il vérifie que
le test évalue la configuration annoncée. Aucune nouvelle modalité des champs encodés
n'est observée en 2024, mais les proportions et relations entre modalités peuvent
évoluer malgré cette stabilité du vocabulaire.

Le test porte sur \textbf{117\,278 usagers}, dont 21\,024 décès ou hospitalisations,
soit 17,93\,\%. Toutes les gravités sont connues dans ce périmètre de 2024.

\section{Des performances proches de celles de 2023}
\begin{table}[htbp]
\centering\small
\renewcommand{\arraystretch}{1.18}
\caption{Scores de validation et de test, configuration inchangée.}\label{tab:ml-test}
\begin{tabularx}{\textwidth}{@{}lLrrrr@{}}
\toprule
Année & Modèle & Log-loss & Brier & AP & ROC-AUC \\
\midrule
2023 & Référence constante & 0,4709 & 0,1474 & 0,1797 & 0,5000 \\
2023 & Régression logistique & 0,4008 & 0,1271 & 0,4085 & 0,7699 \\
2024 & Référence constante & 0,4703 & 0,1471 & 0,1793 & 0,5000 \\
2024 & Régression logistique & 0,4001 & 0,1268 & 0,4090 & 0,7699 \\
\bottomrule
\end{tabularx}
\sourcecalcul{notebook 07 ; apprentissage conservé sur 2019-2022}
\end{table}

La log-loss de la logistique est de 0,4001, contre 0,4703 pour la référence.
La réduction relative est de 14,93\,\%. Le Brier passe de 0,1471 à 0,1268,
soit une réduction de 13,84\,\%. L'AP atteint 0,4090 et la ROC-AUC 0,7699.
Les niveaux restent proches de ceux de 2023.

Ce maintien est un résultat favorable pour le protocole testé~: le gain observé
en validation ne disparaît pas l'année suivante. Il ne garantit pas le comportement
sur d'autres années, une autre population ou des modalités de collecte différentes.
Les écarts de score sont descriptifs~; aucun intervalle supposant les usagers
indépendants n'est utilisé pour revendiquer une significativité.

\section{Pourquoi le seuil 0,5 donne une lecture incomplète}
Pour illustrer les erreurs, une probabilité au moins égale à 0,5 est transformée
en prédiction positive. Le seuil est identique à celui fixé avant le test.
\begin{table}[htbp]
\centering\small
\renewcommand{\arraystretch}{1.18}
\caption{Matrice de confusion sur 2024, seuil illustratif de 0,5.}\label{tab:ml-confusion}
\begin{tabularx}{\textwidth}{@{}Lrr@{}}
\toprule
Issue observée & Prédiction 0 & Prédiction 1 \\
\midrule
Indemne ou blessé léger & 92\,962 & 3\,292 \\
Décès ou hospitalisation & 17\,194 & 3\,830 \\
\textbf{Total} & 110\,156 & 7\,122 \\
\bottomrule
\end{tabularx}
\sourcecalcul{notebook 07 ; classes observées en lignes, classes prédites en colonnes}
\end{table}


Le modèle classe 7\,122 usagers comme positifs. Parmi eux, 3\,830 sont effectivement
décédés ou hospitalisés, soit une précision de 53,78\,\%. Il retrouve cependant
seulement 18,22\,\% des 21\,024 cas positifs~: 17\,194 ne sont pas classés positifs.
Les expressions «~alerte~» ou «~cas manqué~» décrivent ici le tableau de classification,
sans désigner un dispositif de secours.

L'exactitude globale est de 82,53\,\%, contre 82,07\,\% pour la référence qui classe
tout le monde négatif à ce seuil. L'exactitude seule serait donc peu informative~:
elle est élevée même lorsqu'aucun cas positif n'est retrouvé.

\figurerapportML[0.88]{07_02_precision_rappel_test.png}{Courbes précision-rappel de validation et de test pour la configuration inchangée.}{fig:pr-test}{07}
La proximité des courbes résume la stabilité du classement, mais ne fournit pas un
seuil «~correct~». Un autre compromis demanderait un objectif, des coûts d'erreur et
une population d'usage définis, ainsi qu'une évaluation appropriée. Aucun réglage
supplémentaire n'est effectué après lecture de 2024.

\section{La calibration demeure imparfaite}
\figurerapportML[0.98]{07_01_calibration_test.png}{Calibration descriptive sur 2023 et 2024, avec distribution des probabilités du test.}{fig:calibration-test}{07}
Les déformations de la courbe de validation se retrouvent dans le test. Dans
l'intervalle $]0{,}7;0{,}8]$, 1\,030 usagers ont une probabilité moyenne de 74,28\,\%,
alors que 59,90\,\% sont positifs. Le modèle est donc trop assuré dans cette zone.
À l'autre extrémité, les probabilités au plus égales à 0,1 ont une moyenne de
5,22\,\% pour 4,25\,\% observés, sur 42\,619 usagers.

Ces chiffres expliquent pourquoi une ROC-AUC satisfaisante ne suffit pas à interpréter
une probabilité comme une fréquence fiable. Une recalibration pourrait constituer
un travail ultérieur, mais elle ne pourrait pas être choisie puis évaluée sur le
même test 2024 comme si celui-ci était resté indépendant de cette décision.

\section{Examiner les groupes et le rôle des absences}
\begin{table}[htbp]
\centering\small
\renewcommand{\arraystretch}{1.18}
\caption{Diagnostics du test selon la catégorie d'usager.}\label{tab:ml-groupes}
\begin{tabularx}{\textwidth}{@{}Lrrrrr@{}}
\toprule
Catégorie & Usagers & Cas positifs & Observé (\%) & Prédit (\%) & ROC-AUC \\
\midrule
Conducteur & 86\,838 & 14\,608 & 16,82 & 17,04 & 0,7921 \\
Passager & 21\,526 & 3\,514 & 16,32 & 16,76 & 0,6970 \\
Piéton & 8\,914 & 2\,902 & 32,56 & 32,81 & 0,6154 \\
\bottomrule
\end{tabularx}
\sourcecalcul{notebook 07 ; les deux pourcentages sont une proportion observée et une probabilité moyenne}
\end{table}

La logistique améliore la log-loss par rapport à la référence dans les trois
catégories. Le classement reste toutefois moins performant chez les piétons~:
la ROC-AUC y est de 0,6154, contre 0,7921 chez les conducteurs et 0,6970 chez les
passagers. Le graphique des moyennes reste proche des proportions observées.
\figurerapportML[0.88]{07_03_diagnostics_categories_test.png}{Probabilités moyennes et proportions observées par catégorie d'usager en 2024.}{fig:groupes-test}{07}

Les résultats globaux sont majoritairement portés par les conducteurs. Ils ne doivent
pas masquer la difficulté à classer les piétons entre eux. Les scores par tranche
d'âge sont également exportés~; ces diagnostics ne servent pas à sélectionner un
modèle distinct pour chaque groupe.

En 2024, le groupe dont l'âge ou le sexe est absent comprend 2\,424 usagers et un
seul cas positif. Sa très faible log-loss contribue favorablement au score global.
Une ROC-AUC calculée sur un seul cas positif ne constitue pas une preuve solide
de classement dans ce groupe. Le diagnostic de complétude se lit donc avec les
effectifs, et non avec la seule valeur de l'indicateur.

Sur les 114\,854 usagers dont les deux informations sont présentes, le modèle conserve
un avantage~: log-loss de 0,4084 contre 0,4761 pour la référence. Les mêmes prédictions
sont utilisées, sans réentraînement. Cette analyse montre que le gain n'est pas limité
aux absences d'âge ou de sexe~; elle ne démontre ni la cause de ces absences, ni la
robustesse à une modification des pratiques de saisie.

\section{Ce que cette expérience permet de conclure}
L'étude montre qu'un modèle simple, ajusté sur les années antérieures, estime les
probabilités mieux qu'une fréquence constante sur 2023 et 2024. Elle donne aussi
des diagnostics concrets~: faibles différences entre plusieurs valeurs de $C$,
surconfiance dans les intervalles élevés, rappel limité au seuil illustratif et
performances inégales selon les usagers.

Les limites principales restent celles de la question et des données. BAAC contient
les accidents corporels enregistrés, avec une cible d'hospitalisation sensible à la
collecte. Les usagers d'un accident sont dépendants. La représentation en tranches
et le modèle additif ne capturent pas toutes les relations possibles. Les variables
disponibles dans le fichier final ne définissent pas un instant de prédiction réel.
Enfin, une seule année de test ne suffit pas à établir une stabilité générale.

La valeur ajoutée du volet ML est donc une évaluation probabiliste contrôlée et
une lecture critique des résultats. Une aide à la décision demanderait une autre
étape de cadrage, des informations disponibles au moment voulu et une validation
de l'usage. Elle n'est pas démontrée par ces scores.

\FloatBarrier
\bilan{Le gain par rapport à la référence se maintient en 2024, avec une log-loss
proche de 0,400 et une ROC-AUC proche de 0,770. La conclusion reste mesurée~:
le modèle contient un signal prédictif, mais ses probabilités et ses erreurs ne
justifient pas à elles seules un usage opérationnel.}
